\documentclass[11pt]{article}
\usepackage[utf8]{inputenc}
\usepackage[T1]{fontenc}
\usepackage[margin=1in]{geometry}
\usepackage{hyperref}
\usepackage{enumitem}
\title{Grade 2 - Day 1 Community Project}
\author{Piter Garcia}
\date{March 20, 2026}
\begin{document}
\maketitle
\section*{Quick Scan}
\begin{itemize}[leftmargin=*]
\item Grade: Grade 2
\item Workbook sessions: Session 32
\item Class blocks: Day 1 10:45-11:35 Callon, Day 1 2:15-3:05 Henchen, Day 4 10:45-11:35 Barthelman
\item Reference file: grade-2-day-1-community-project.bib
\end{itemize}
\section*{School Planning Goals and Inclusion}
\begin{itemize}[leftmargin=*]
\item What is expected: I can collaboratively plan and design urban, suburban, and rural communities using models and drawings. | | I can create a model of an object or computational process in order to identify patterns and essential elements of the object or process. | | I can Identify multiple ways that the same problem could be decomposed into smaller steps. | | I can develop and document a plan that outlines specific steps taken to complete a project.
\item What we achieved: No direct transcript match is attached yet. Treat this lesson as workbook-backed and enrich it when a matching classroom transcript is available.
\item What needs improvement: stronger proof, cleaner assessment notes, and tighter lesson artifacts.
\item How to improve: add transcript proof, one saved artifact, and clearer success criteria.
\item Inclusion focus: use visual modeling, chunked directions, predictable routines, and flexible participation roles.
\end{itemize}
\section*{Official References}
\begin{itemize}[leftmargin=*]
\item \url{https://csteachers.org/k12standards/interactive/}

\end{itemize}
\end{document}
